\section{Processus d'aménagement des plots : étapes, rôle et justification des apports}
\subsection{Les étapes de l'aménagement}
\subsubsection{Désherbage et Isolement des herbes}
\begin{itemize}[label=$\bullet$]
   	 \item Le désherbage consiste à nettoyer les herbes sur une parcelle avant la préparation du sol. Cette action permet d’assainir la surface, facilitant les travaux de labour.
	\end{itemize}

\begin{itemize}[label=$\bullet$]
   	 \item L’isolement des herbes consiste à regrouper et retirer les herbes  hors de la zone de travail afin de maintenir la zone propre.
	\end{itemize}

    \begin{figure}[H]
    \centering
    \imgstd{images/08_desherbage_isolement_herbes.jpg}
    \caption{Processus de désherbage et d’isolement des herbes : début et fin. }
    \label{fig:desherbage2}
\end{figure}

\subsubsection{Le labour et la délimitation des plots}
Après un labour profond sur les trois plots constituant la zone test, nous avons ajouté du biochar, de la PCA, de l'EGB et du compost que nous avons soigneusement mélangés avant d'aplanir et de fixer les ados. En revanche, dans la zone témoin, nous avons simplement délimité et fixé les ados après un labour approfondi.

\begin{itemize}[label=$\bullet$]
\item Le labour est une pratique agricole qui consiste à retourner la couche superficielle de la terre. Cette opération permet d'aérer et d’ameublir le sol pour préparer les cultures à venir.

\item Le travail en profondeur mené lors de la réalisation des plots a pour but d'ameublir le sol, d'améliorer son aération et sa teneur en eau, tout en stimulant la croissance des racines des cultures à venir et l'intégration des amendements ajoutés.

\item La confection de plots avec des ados a contribué à délimiter les espaces cultivés et les voies de passage, rendant par conséquent les opérations et la supervision de l'irrigation plus aisées.
\end{itemize}

\begin{figure}[H]
    \centering
\imgstd{images/11_mise_en_place_plot_irrigation_fixer_ados.jpg}
    \caption{Mise en place des plots et irrigation pour la fixation des ados. }
    \label{fig:plotAdos2}
\end{figure}

\subsection{Rôle et justification du choix des apports(mécanisme attendu)}
Les éléments apportés sont généralement des éléments de correction du sol:
\subsubsection{Le biochar}
Le biochar possède une surface poreuse et aromatique riche en groupes fonctionnels chargés 
négativement (carboxyles \ce{-COOH}, hydroxyles \ce{-OH}, phénols \ce{Ar-OH}), ce qui lui donne 
une capacité d'échange cationique (CEC) élevée \cite{chandrasekaran_comprehensive_2026}. Cette CEC lui permet d'attirer
et de retenir des ions chargés positivement appelés cations du sol, dont le \ce{Na+} (sodium), qui est
principalement responsable de la salinité. Le mécanisme précis pour lutter contre la salinité 
repose sur un échange préférentiel : les cations polyvalents présents dans le biochar, notamment 
le calcium (\ce{Ca^2+}) et le magnésium (\ce{Mg^2+}), peuvent remplacer le \ce{Na+} sur les sites 
d'échange, permettant ensuite à l'excès de sodium libéré d'être facilement lessivé grâce à 
l'irrigation ou au ruissellement \cite{yuan_biochar_2023}. De plus, le biochar favorise l'absorption du potassium par la plante
tout en diminuant celle du sodium, participant à l'équilibre ionique et à la régulation osmotique de la plante \cite{chandrasekaran_comprehensive_2026,zubairu_overview_2025}.

    \begin{figure}[H]
    \centering
\imgstd{images/12_biochar.jpg}
    \caption{Poudre de charbon de bois (dit Biochar en agriculture). }
    \label{fig:biochar2}
\end{figure}

\subsubsection{La poudre de coque d’arachide}
La CA est principalement composée de 45~\% de cellulose, de 6~\% d’hémicellulose et de
36~\% de lignine, ce qui en fait un bon amendement pour le sol; sa teneur élevée en lignine
entraîne cependant une décomposition lente dans l’environnement naturel \cite{sawargaonkar_valorization_2024}. Grâce à sa
décomposition lente, elle représente une source durable de matière organique qui améliore
\textbf{l’aération et la structure} du sol sur le long terme, contrairement à un apport nutritif rapide
comme le compost. Une étude menée spécifiquement sur des sols salins a montré que
l’amendement avec des CA (en quantités de 4, 6 ou 8 tonnes par hectare) améliore les
propriétés du sol et la croissance de jeunes plants d’espèces ligneuses cultivées sur des
sols salins \cite{fall_effect_2018}. 

NB : \textbf{Les coquilles d’arachide représentent 30~\% du poids des arachides.}

\begin{figure}[H]
    \centering
\imgstd{images/13_coque_arrachide.jpg}
    \caption{Poudre de coque d’arachide. }
    \label{fig:arrachide2}
\end{figure}

\subsubsection{Le compost}
Le compost améliore les propriétés physiques, chimiques et biologiques du sol : il augmente la capacité de rétention en eau du sol, 
la disponibilité des nutriments nécessaires ($N$, $P$, $K$) et renforce la résistance des cultures face aux maladies \cite{goldan_assessment_2023}. 
Contrairement à un engrais minéral à libération immédiate, le compost libère ses éléments nutritifs de façon progressive, assurant ainsi un apport constant au sol et réduisant le besoin de fertilisations répétées.

En ce qui concerne plus spécifiquement la salinité, une étude portant sur l’irrigation
avec une eau saumâtre a montré qu’un compost enrichi biologiquement a contribué à
la réduction de l’accumulation de sodium de 60~\% et celle du chlore de 70~\%, tout en
améliorant la rétention d’eau et la structure du sol \cite{suvendran_soil_2025}.

\begin{figure}[H]
    \centering
\imgstd{images/14_compost.jpg}
    \caption{Compost (Matière organique). }
    \label{fig:compost2}
\end{figure}

\subsubsection{Le biofertilisant anti-sel ou engrais granulé bio}
Contrairement aux trois amendements précédents, dont l’action est essentiellement physicochimique, le EGB agit biologiquement, via des micro-organismes rhizosphériques. 
Les rhizobactéries halotolérantes favorisant la croissance des plantes (HT-PGPR, Halotolerant Plant Growth-Promoting Rhizobacteria) agissent par plusieurs mécanismes combinés : 
la production de métabolites secondaires et de sidérophores, la synthèse de phytohormones et d’enzymes, le maintien de l’équilibre ionique, et l’amélioration de la disponibilité en nutriments \cite{ramasamy_coping_2023}.

Le mécanisme d’équilibrage ionique est particulièrement efficace pour la lutte contre la
salinité : sous stress salin, les plantes réduisent leur évapotranspiration en renforçant l’absorption 
de potassium (K+) et en diminuant celle du sodium (Na+), ce qui atténue le stress osmotique et ionique. 
Ce processus est facilité par les bactéries PGPR qui augmentent l’absorption des nutriments bénéfiques par la plante \cite{gao_progress_2022}.

\begin{figure}[H]
    \centering
\imgstd{images/15_anti_sel.jpg}
    \caption{Engrais granulé bio. }
    \label{fig:EGB2}
\end{figure}
    

    

  




 





 

